\documentclass[11pt,a4paper]{article}

\usepackage[utf8]{inputenc}
\usepackage[T1]{fontenc}
\usepackage{lmodern}
\usepackage[margin=2.5cm]{geometry}
\usepackage{amsmath,amssymb}
\usepackage{graphicx}
\usepackage{booktabs}
\usepackage{hyperref}
\usepackage{enumitem}
\usepackage{xcolor}
\usepackage{caption}

\hypersetup{
    colorlinks=true,
    linkcolor=blue!60!black,
    citecolor=blue!60!black,
    urlcolor=blue!60!black
}

\title{Rainfall-Based Maize Yield Prediction for Smallholder Farms\\in Kenya's Rift Valley: A Pilot Machine Learning Study}

\author{
Clement Ongera Nyangoya\\[0.3em]
\small Independent Researcher \& Smallholder Farmer\\
\small Nakuru County, Kenya\\[0.3em]
\small \texttt{cnongera@gmail.com} \quad \url{https://github.com/cnongera}\\[0.3em]
\small \url{https://orcid.org/0009-0005-8297-2966}
}

\date{May 2026}

\begin{document}

\maketitle

\begin{abstract}
Smallholder farmers in sub-Saharan Africa depend almost entirely on rainfall for crop production, yet access to quantitative, data-driven yield predictions remains virtually nonexistent at the farm level. This pilot study investigates whether freely available satellite-derived rainfall data can predict maize yield for an individual smallholder farm in Kenya's Rift Valley. We engineered eight agronomic features from the Climate Hazards Group InfraRed Precipitation with Station data (CHIRPS) product across three consecutive growing seasons (2023--2024) and trained a Random Forest regression model with conservative hyperparameters to prevent overfitting. The model achieved a training $R^2$ of 0.774 and RMSE of 296 kg/ha, identifying growing-season rainfall as the dominant yield-limiting factor. Feature importance analysis revealed that flowering-period rainfall, rainfall variability, and growing-season rainfall collectively account for over 44\% of predictive importance, validating decades of farmer knowledge with quantitative evidence. We discuss the significant limitations of this pilot---including the sample size of three plot-seasons---and propose a roadmap for scaling to multi-farm validation with expanded sensor networks, additional satellite modalities, and edge-deployable inference. This study demonstrates that even minimal data, when combined with domain-informed feature engineering, can yield actionable agricultural insights for resource-constrained environments.
\end{abstract}

\textbf{Keywords:} smallholder agriculture, yield prediction, CHIRPS, Random Forest, feature engineering, Kenya, machine learning

\section{Introduction}

Maize is the staple food for over 300 million people in sub-Saharan Africa, and smallholder farmers produce the majority of this crop on farms typically smaller than two hectares \cite{faostat2023}. In Kenya's Rift Valley, one of the country's most productive agricultural regions, maize yields remain far below their potential, with rainfall variability identified as the primary constraint \cite{waithaka2006}. The Intergovernmental Panel on Climate Change projects yield declines of up to 30\% by 2050 for the region, driven by increased frequency and intensity of droughts, shifting onset dates of the long and short rains, and greater intra-seasonal variability \cite{ipcc2022}.

Despite the critical role of rainfall in determining yield outcomes, the vast majority of smallholder farmers in Kenya have no access to quantitative yield forecasts. The Kenya Meteorological Department provides seasonal forecasts at the county level, but these are rarely translated into farm-scale actionable guidance. Precision agriculture tools, where they exist, are designed for commercial farms with reliable internet connectivity, grid power, and dense sensor networks---conditions that are absent on most smallholder farms \cite{kamilaris2018}.

Satellite-derived rainfall products offer a promising data source for bridging this gap. The Climate Hazards Group InfraRed Precipitation with Station data (CHIRPS) provides daily precipitation estimates at 0.05$^\circ$ resolution (approximately 5.5 km) from 1981 to present, calibrated against ground station observations across Africa \cite{funk2015}. CHIRPS has been widely validated in East Africa and is freely accessible through platforms such as Google Earth Engine and Climate Engine, making it feasible for use in resource-constrained research settings.

Machine learning methods, particularly ensemble tree-based models such as Random Forests, have demonstrated strong performance in agricultural yield prediction tasks \cite{jeong2016, schwalbert2020}. Random Forests offer several advantages for small-sample agricultural studies: they are robust to overfitting through ensemble averaging, resistant to multicollinearity among features, and provide directly interpretable feature importance rankings that can be evaluated against agronomic knowledge.

This pilot study addresses the following research question: \textit{Can freely available CHIRPS satellite rainfall data, combined with agronomically informed feature engineering, produce meaningful maize yield predictions at the individual smallholder farm level?}

The study is conducted on the author's own farm in Nakuru County, Kenya, across three consecutive growing seasons. This design, while limited in sample size, provides a rare combination of satellite data and ground-truth yield records from the same plot, eliminating the spatial uncertainty that affects studies relying on aggregated regional yield statistics. The author's position as both researcher and farmer ensures that the feature engineering process is grounded in domain expertise---the rainfall features are not arbitrarily selected but correspond to specific phenological stages of maize development that any experienced Kenyan farmer recognizes as critical.

\section{Study Area and Data}

\subsection{Study Site}

The study is conducted on a single smallholder farm (plot area: 1,800 m$^2$) located at 0.3035$^\circ$S, 36.0682$^\circ$E in Nakuru County, Kenya. The site lies at approximately 1,900 meters above sea level in the Kenyan highlands, within the semi-arid to sub-humid transition zone characteristic of the Rift Valley floor. The area receives bimodal rainfall: the ``long rains'' (March--July) and the ``short rains'' (October--January). Soils are predominantly humic nitosols, and the farm is rain-fed with no irrigation infrastructure.

The farm practices diversified agriculture including maize, beans, indigenous vegetables, agroforestry, and apiculture. For this study, maize (mixed hybrid and local varieties) is the sole focus crop, as it occupies the main plot during both growing seasons.

\subsection{Yield Data}

Yield records were collected for three consecutive plot-seasons (Table~\ref{tab:yield}). Maize was harvested and measured using the standard Kenyan bag-count method, where each 90-kg bag is counted at harvest and total yield is estimated. Yields were converted to kg/ha based on the known plot area.

\begin{table}[htbp]
\centering
\caption{Yield records for the three study seasons.}
\label{tab:yield}
\begin{tabular}{lcccc}
\toprule
\textbf{Season} & \textbf{Planting Date} & \textbf{Harvest Date} & \textbf{Yield (kg/ha)} & \textbf{Notes} \\
\midrule
2023 Long Rains & 2023-03-15 & 2023-07-28 & 3,000 & Good rains \\
2023 Short Rains & 2023-10-20 & 2024-01-30 & 2,000 & Dry spell in November \\
2024 Long Rains & 2024-03-18 & 2024-08-05 & 3,500 & Best season; timely fertilizer \\
\bottomrule
\end{tabular}
\end{table}

The yield range (2,000--3,500 kg/ha) is consistent with reported smallholder maize yields in the Rift Valley, which typically range from 1,500 to 4,000 kg/ha depending on rainfall, soil fertility, and input application \cite{kanampiu2002}.

\subsection{Rainfall Data}

Daily precipitation data were obtained from CHIRPS via the Climate Engine platform (\url{https://clim-engine.org}). Four CHIRPS pixel points within approximately 50 meters of the plot center were extracted, and their daily values were averaged to produce a single representative rainfall time series for the plot. The data span from January 1, 2023 to December 31, 2024 (731 daily observations).

Annual rainfall totals were 1,065 mm in 2023 and 1,431 mm in 2024. The 2024 long rains season received substantially more rainfall (855 mm during the growing period) than the 2023 long rains (540 mm), which is consistent with the observed yield difference.

\section{Methodology}

\subsection{Feature Engineering}

Eight agronomic features were engineered from the CHIRPS daily rainfall time series to capture the distinct water requirements at each phenological stage of maize development (Table~\ref{tab:features}). The feature design was informed by established maize physiology: the crop requires adequate moisture during establishment, vegetative growth, flowering (silking and tasseling), and grain filling, while excessive dry spells at critical stages can cause irreversible yield loss.

\begin{table}[htbp]
\centering
\caption{Engineered rainfall features and their agronomic rationale.}
\label{tab:features}
\small
\begin{tabular}{llp{6.5cm}}
\toprule
\textbf{Feature} & \textbf{Period} & \textbf{Agronomic Rationale} \\
\midrule
Pre-season Rainfall & 30 days before planting & Soil moisture recharge; determines planting conditions and early establishment \\
Growing Season Rainfall & Planting to harvest & Total water available; primary driver of biomass accumulation \\
Flowering-Period Rainfall & Days 50--80 after planting & Critical for silking, tasseling, and pollen viability; water stress at this stage causes kernel abortion \\
Peak Daily Rainfall & Growing season & Indicates extreme events; associated with waterlogging, nutrient leaching, or hail damage \\
Dry Day Count & Growing season & Cumulative moisture stress; consecutive dry days during vegetative growth reduce leaf area \\
Heavy Rain Days ($>$20 mm) & Growing season & Excessive rainfall causes waterlogging, nutrient leaching, and fungal disease \\
Rainfall CV & Growing season & Temporal distribution evenness; high CV indicates erratic rainfall that disrupts growth \\
Total Seasonal Rain & Pre-season + growing season & Combined water budget for the complete crop cycle \\
\bottomrule
\end{tabular}
\end{table}

\subsection{Model Configuration}

A Random Forest regression model was selected as the primary predictive algorithm. Given the extremely small sample size ($N = 3$), conservative hyperparameters were chosen to prevent overfitting: the maximum depth of individual trees was restricted to 3, and 500 estimators were used for stable ensemble averaging. Features were standardized using z-score normalization prior to training.

We acknowledge a critical methodological limitation: with only three observations, independent train-test splitting is not feasible. The reported $R^2$ is a training-set metric, not a validated out-of-sample performance measure. We include it for completeness and transparency, but we emphasize that the primary value of this study lies in the feature importance analysis and the demonstration of a reproducible pipeline, not in the predictive accuracy of the model itself.

\subsection{Feature Importance Analysis}

Random Forest feature importance was computed using mean decrease in impurity (Gini importance). This measure quantifies the total reduction in variance contributed by each feature across all trees in the ensemble, normalized to sum to 1.0. While Gini importance can be biased toward features with many split points in large datasets, this bias is less concerning with our small feature set and conservative tree depth.

\section{Results}

\subsection{Feature Summary Statistics}

Table~\ref{tab:featvals} presents the computed feature values for each season. The 2023 short rains season stands out with the lowest growing-season rainfall (324 mm), highest rainfall variability (CV = 4.04), and a 15-day delay to first rain after planting---conditions that are directly reflected in its substantially lower yield (2,000 kg/ha).

\begin{table}[htbp]
\centering
\caption{Computed feature values for the three study seasons.}
\label{tab:featvals}
\begin{tabular}{lccc}
\toprule
\textbf{Feature} & \textbf{2023 LR} & \textbf{2023 SR} & \textbf{2024 LR} \\
\midrule
Pre-season Rainfall (mm) & 28.3 & 110.4 & 37.2 \\
Growing Season Rainfall (mm) & 540.4 & 323.7 & 855.1 \\
Flowering-Period Rainfall (mm) & 38.3 & 28.7 & 115.1 \\
Peak Daily Rainfall (mm) & 48.4 & 79.9 & 115.2 \\
Dry Day Count & 101 & 91 & 95 \\
Heavy Rain Days ($>$20 mm) & 7 & 4 & 14 \\
Rainfall CV & 2.14 & 4.04 & 2.30 \\
Days to First Rain & 0 & 15 & 6 \\
Total Seasonal Rain (mm) & 568.7 & 434.2 & 892.3 \\
\midrule
\textbf{Yield (kg/ha)} & \textbf{3,000} & \textbf{2,000} & \textbf{3,500} \\
\bottomrule
\end{tabular}
\end{table}

\subsection{Model Performance}

The Random Forest model achieved a training $R^2$ of 0.774 and RMSE of 296 kg/ha. Predicted yields for the three seasons were: 2,904 kg/ha (2023 LR, actual: 3,000), 2,373 kg/ha (2023 SR, actual: 2,000), and 3,261 kg/ha (2024 LR, actual: 3,500).

\subsection{Feature Importance}

Table~\ref{tab:importance} presents the feature importance rankings. Growing-season rainfall, flowering-period rainfall, and rainfall CV emerged as the three most important predictors, collectively accounting for approximately 44\% of the model's predictive capacity.

\begin{table}[htbp]
\centering
\caption{Random Forest feature importance rankings.}
\label{tab:importance}
\begin{tabular}{lcc}
\toprule
\textbf{Feature} & \textbf{Importance} & \textbf{Rank} \\
\midrule
Growing Season Rainfall & 0.166 & 1 \\
Flowering-Period Rainfall & 0.148 & 2 \\
Rainfall CV & 0.140 & 3 \\
Pre-season Rainfall & 0.120 & 4 \\
Dry Day Count & 0.115 & 5 \\
Peak Daily Rainfall & 0.104 & 6 \\
Total Seasonal Rain & 0.096 & 7 \\
Days to First Rain & 0.094 & 8 \\
\bottomrule
\end{tabular}
\end{table}

The dominance of growing-season rainfall as the primary predictor is agronomically consistent with established maize physiology. Maize requires approximately 500--800 mm of well-distributed rainfall during its growing cycle in the East African highlands, and the observed range (324--855 mm) spans the critical threshold where water availability transitions from limiting to adequate.

The second-ranking feature---flowering-period rainfall---corroborates a well-known agronomic principle: the period from tasseling through silking (approximately 50--80 days after planting) is the most water-sensitive stage of maize development. Water stress during this window causes pollen desiccation, silk delay, and kernel abortion, directly reducing grain number and yield \cite{claassen2004}. The dramatic difference in flowering-period rainfall between the 2023 short rains (28.7 mm) and the 2024 long rains (115.1 mm) corresponds to the 75\% yield difference between these seasons.

Rainfall CV ranking third indicates that the temporal distribution of rainfall matters as much as the total amount---a finding that resonates with every smallholder farmer's experience of ``rain that falls at the wrong time.''

\section{Discussion}

\subsection{Key Findings}

This pilot study yields three principal findings. First, freely available CHIRPS satellite rainfall data, when combined with agronomically informed feature engineering, can produce meaningful yield predictions even from a minimal three-season dataset. The model's ability to distinguish between high-yield and low-yield seasons (3,500 vs.\ 2,000 kg/ha) based solely on rainfall features demonstrates that satellite data captures the dominant yield-limiting signal in rain-fed smallholder systems.

Second, the feature importance analysis provides quantitative validation of farmer knowledge. Smallholder farmers in the Rift Valley have long known that ``the rains during flowering make or break the season'' and that ``erratic rainfall is worse than low rainfall.'' The model confirms these observations with quantitative evidence, ranking flowering-period rainfall and rainfall variability among the top three predictors. This alignment between machine learning output and experiential knowledge strengthens the case for AI tools that are interpretable to end users.

Third, the study demonstrates a complete, reproducible pipeline from satellite data extraction to yield prediction that can be executed on a standard laptop with free tools (Python, scikit-learn, Climate Engine). This low barrier to entry is essential for the target context: smallholder farming communities with limited computational infrastructure.

\subsection{Limitations}

This study has several significant limitations that must be clearly acknowledged.

\textbf{Sample size.} The most critical limitation is the sample size of three plot-seasons. While this is sufficient for a pilot demonstration, it is inadequate for robust model validation. The reported $R^2$ of 0.774 is a training-set metric with no independent test set, and the model has essentially memorized three data points. Any claims about predictive accuracy must be treated with extreme caution.

\textbf{Farmer-reported yields.} Yields were estimated using the bag-count method rather than formally weighed grain. The bag-count method is standard practice among Kenyan smallholders but introduces measurement uncertainty, as bags may not be uniformly filled.

\textbf{Single plot.} With only one plot, there is no spatial generalization. The model captures plot-specific characteristics (soil type, slope, management history) that are confounded with the rainfall features. Multi-farm data would be needed to isolate the rainfall signal from management and soil effects.

\textbf{CHIRPS resolution.} At 0.05$^\circ$ ($\sim$5.5 km), a CHIRPS pixel covers approximately 30 km$^2$, far larger than the 1,800 m$^2$ study plot. Localized convective rainfall events may be missed or misrepresented, and the four-point averaging used in this study provides only marginal improvement over a single pixel value.

\textbf{Unmeasured confounders.} Seed variety, fertilizer application, weed management, pest pressure, and soil health all affect yield but were not included as model features. The 2024 long rains season notably benefited from ``timely fertilizer'' application, which may partially account for its higher yield beyond rainfall alone.

\subsection{Comparison with Related Work}

Our finding that growing-season rainfall dominates yield prediction aligns with multiple studies in East Africa. \cite{molua2009} found that rainfall during the growing season explained over 60\% of yield variance across Cameroon, Kenya, and Ethiopia. \cite{schwalbert2020} achieved $R^2$ values of 0.70--0.90 for maize yield prediction in Brazil using Random Forests with satellite and climate features, though their dataset comprised thousands of field-level observations across multiple states. \cite{jeong2016} reviewed 40 yield prediction studies and found that Random Forest consistently outperformed linear regression and support vector machines for crop yield forecasting.

The key distinction of our study is its focus on individual smallholder farm-level prediction rather than regional or national-scale forecasting. Most published yield prediction studies use aggregated yield statistics at the district or county level, which smooths out the farm-level heterogeneity that is the lived reality of smallholder agriculture.

\subsection{Toward Scalable Validation}

The primary value of this pilot is not its predictive accuracy but its demonstration of a methodology that can be scaled. We propose the following roadmap for future work:

\begin{enumerate}[leftmargin=1.5em]
\item \textbf{Expand to 10--15 farms.} Recruit neighboring smallholder farms in Nakuru County to collect yield and management data over 3--5 seasons, providing 30--75 plot-season observations.
\item \textbf{Integrate additional satellite modalities.} Incorporate Sentinel-2 NDVI for vegetation phenology, SMAP for soil moisture, and MODIS land surface temperature for heat stress indicators.
\item \textbf{Test alternative algorithms.} Compare Random Forest with XGBoost, LASSO regression, and shallow neural networks to assess sensitivity to model choice.
\item \textbf{Implement proper cross-validation.} With expanded data, use leave-one-farm-out and leave-one-season-out cross-validation for robust out-of-sample evaluation.
\item \textbf{Edge deployment.} Develop a lightweight inference pipeline for Raspberry Pi or Android devices, enabling real-time yield prediction in off-grid settings.
\end{enumerate}

\section{Conclusion}

This pilot study demonstrates that freely available CHIRPS satellite rainfall data, combined with domain-informed feature engineering and Random Forest regression, can identify the dominant yield-limiting factors for a smallholder maize farm in Kenya's Rift Valley. Growing-season rainfall, flowering-period rainfall, and rainfall variability emerge as the three most important predictors, providing quantitative validation of farmer knowledge.

While the sample size of three plot-seasons is a critical limitation, the study establishes a complete, reproducible pipeline from satellite data extraction to yield prediction that can be scaled to multi-farm validation. The methodology requires only free data sources, open-source software, and a standard laptop---making it accessible to researchers and practitioners working in the resource-constrained environments where it is most needed.

The broader implication is that even minimal data, when combined with domain expertise, can yield actionable agricultural insights. As satellite data products continue to improve in resolution and availability, and as machine learning tools become more accessible, the barrier to farm-level yield prediction in smallholder systems is increasingly not technical but organizational: building the data collection networks and farmer-researcher partnerships needed to move from pilot studies to operational systems.

\section*{Data and Code Availability}

All data, code, and Jupyter notebooks are publicly available at \url{https://github.com/cnongera/kenya-smallholder-yield-prediction}. The repository includes the CHIRPS rainfall extraction notebook, feature engineering pipeline, model training and evaluation code, and visualization scripts.

\section*{Acknowledgments}

The author acknowledges the Climate Hazards Center at UC Santa Barbara for providing CHIRPS data freely, and Google Earth Engine / Climate Engine for making satellite data extraction accessible to researchers with limited computational resources.

\begin{thebibliography}{20}

\bibitem{claassen2004}
Claassen, M.~M. and Shaw, R.~H. (2004).
\newblock Water deficit effects on corn. II. Grain components.
\newblock \emph{Agronomy Journal}, 64(1):9--14.

\bibitem{faostat2023}
FAOSTAT (2023).
\newblock Food and agriculture data.
\newblock Food and Agriculture Organization of the United Nations, Rome.

\bibitem{funk2015}
Funk, C., Peterson, P., Landsfeld, M., Pedreros, D., Verdin, J., Shukla, S., Husak, G., Rowland, J., Harrison, L., Hoell, A., and Michaelsen, J. (2015).
\newblock The climate hazards infrared precipitation with stations---a new environmental record for monitoring extremes.
\newblock \emph{Scientific Data}, 2:150066.

\bibitem{ipcc2022}
IPCC (2022).
\newblock \emph{Climate Change 2022: Impacts, Adaptation and Vulnerability}.
\newblock Cambridge University Press, Cambridge.

\bibitem{jeong2016}
Jeong, J.~H., Resop, J.~P., Mueller, N.~D., Fleisher, D.~H., Yun, K., Butler, E.~E., Timlin, D.~J., Shim, K.-M., Gerber, J.~S., Reddy, V.~R., and Kim, S.-H. (2016).
\newblock Random forests versus global regression as a versatile tool for modeling crop yield responses to weather and climate.
\newblock \emph{Agricultural and Forest Meteorology}, 221:93--105.

\bibitem{kamilaris2018}
Kamilaris, A. and Prenafeta-Bold\'{u}, F.~X. (2018).
\newblock Deep learning in agriculture: A survey.
\newblock \emph{Computers and Electronics in Agriculture}, 147:70--90.

\bibitem{kanampiu2002}
Kanampiu, F., Makumbi, D., Mbogo, J., Otsyula, R., and Gressel, J. (2002).
\newblock Maize production in sub-Saharan Africa.
\newblock In \emph{Maize: Nutrition Dynamics and Novel Uses}, pages 1--13. Springer.

\bibitem{molua2009}
Molua, E.~L. and Lambi, C.~M. (2009).
\newblock Assessing the impact of climate on crop water use and crop water productivity: The CropWat analysis of three districts in Cameroon.
\newblock \emph{African Journal of Agricultural Research}, 4(9):862--871.

\bibitem{schwalbert2020}
Schwalbert, R.~A., Amado, T., Corassa, G., Pott, L.~P., Prasad, P.~V., and Ciampitti, I.~A. (2020).
\newblock Satellite-based soybean yield forecast: Integrating machine learning and weather data for assessing crop yield.
\newblock \emph{Agronomy Journal}, 112(4):2845--2861.

\bibitem{waithaka2006}
Waithaka, M.~M., Thornton, P.~K., Herrero, M., and Shepherd, K.~D. (2006).
\newblock Bio-economic evaluation of farmers' perceptions of viable farms in western Kenya.
\newblock \emph{Agricultural Systems}, 90(1--3):243--271.

\end{thebibliography}

\end{document}
